% page 148 du fac-simile (image p-147 du scan)
% bloc 160.0 x 237.7 mm ; marge gauche 8.0 mm
\pgc{-12.700mm}{0.000mm}{
\l{\cel{3}138.}
\l{}
\l{}
\l{\sou{eparko}. I. (\sou{che la Romani antiqua}). Guvernisto di provinco}
\l{\cel{3}en lando Greka. - II. (\sou{lor la imperio Greka}). Prefekto}
\l{\cel{3}di Konstantinoplo. - III.(en\cel{1}\sou{Grekia}).\cel{2}Subprefekto.- D-}
\l{EFIRS.}
\l{}
\l{\sou{eperlano}. (\sou{zool.}) Mikra marfisho, briloze perlomatrea, di}
\l{\cel{3}qua la karno esas delikata. - L.\cel{1}\sou{Osmerus eperlanus.-}.FIL.}
\l{}
\l{\sou{epicena}. Qua indikas sen-distinge la ento maskula,e la ento}
\l{\cel{3}femina qua korespondas. - DEFIS.}
\l{}
\l{\sou{epiciklo}. (\sou{astron.)}\cel{2}En laastronomio olima : mikra cirklo}
\l{\cel{3}qua (lore-konjekte) parkuris la periferio di cirklo plu}
\l{\cel{3}granda, tale justifikante la neregulalaji videbla quin onu}
\l{\cel{3}deskovris en la movi di la astri. - DEFIRS.}
\l{}
\l{\sou{epicikloido}. (\sou{geom}.) Kurvo produktata da la rotaco di punto}
\l{\cel{3}selektita sur kurvo movanta qua rulas sen glito cirkum}
\l{\cel{3}kurvo fixa. - DEFIRS.}
\l{}
\l{\sou{epidemio}. Morbo qua atakas, en un loko, en la sama instanto,}
\l{\cel{3}granda quanto de personi, ed aspektas kom dependanta de}
\l{\cel{3}kauzo qua efikas generale. - DEFIRS.}
\l{}
\l{\sou{epidermo}. (\sou{anat.}) Strato surfacala di la pelo, membrano dina}
\l{\cel{3}e diafana qua kovras la dermo. - DEFIRS.}
\l{}
\l{\sou{epifanio}. (\sou{religio kristana}). I. Aparo di Yesu a la Reji Maga}
\l{\cel{3}qui venis adorar lu. - II. Festo-dio Ekleziala (katolika)}
\l{\cel{3}en la dio sisesma di januaro, por celebrar ta adoro.-DEFIRS.}
\l{}
\l{\sou{epifita}. Qua vivas sur altra planti, sen tirar de li sua}
\l{\cel{3}nutrivi. - DE.}
\l{}
\l{\sou{epifizo}. (\sou{anat.}) Saliajo ostatra qua, ligita - che la infanto}
\l{\cel{3}- a la osto per kartilago, divenas pose apofizo. - DE.}
\l{}
\l{\sou{epigastro}. (\sou{anat.)}\cel{2}Parto supera di la abdomino, situita an-}
\l{\cel{3}super la umbiliko. - DEF.}
\l{}
\l{\sou{epigenezo}. Doktrino segun qua onu konceptas la formaceso di la}
\l{\cel{3}korpi vivanta per la adjunto sucedanta di parti qui ne pre-}
\l{\cel{3}existis en la jermo.}
\l{}
\l{\sou{epigloto}. (\sou{anat.}) Valveto ye la parto supra di laringo, di}
\l{\cel{3}qua lu klozas (per sua abaseso) la orifico supra lor la en-}
\l{\cel{3}gluto, por impedar la eniro di la nutrivi aden la tubi}
\l{\cel{3}respirala. - DEF.}
\l{}
\l{\sou{epigono}. I. (\sou{mitol}.) Singla de la heroi di la milito duesma di}
\l{\cel{3}Tebo. - II. (\sou{metaf}.) La homo qua apartenas a la generaciono}
\l{\cel{3}dato-duesma. - DEF.}
}
